\section*{Capítulo \ref{ch:b1:radiometric-methods} --- Métodos de datación radiométrica}
% values: calc/ch25.py

\begin{solution}{exo:b1:radiometric-methods:1}
(a)~$D^*/P = 2^n - 1 = 3$ da $n = 2$: dos semividas. (b)~Rubidio-87:
$e^{\lambda t} - 1 = e^{\num{1.3972e-11} \times \num{3.10e8}} - 1 = 0.004341$;
uranio-235: $e^{\num{9.8485e-10} \times \num{3.10e8}} - 1 = 0.3570$. El
cociente del uranio es unas 80 veces mayor: un error absoluto dado de medida es una parte mucho menor de él. (c)~$2^{-n} < 0.01$ cuando $n > \log_2 100 = 6.64$: tras
unas siete semividas.
\end{solution}

\begin{solution}{exo:b1:radiometric-methods:2}
(a)~Radiocarbono sobre el carbón vegetal. (b)~Uranio--plomo sobre el circón, o
argón--argón (potasio--argón) sobre la sanidina, un \omterm{def:b1:rock-forming-minerals:mineral}{mineral} rico en potasio.
(c)~Uranio--plomo sobre circón, que sobrevive a una historia larga; el rubidio--estroncio sobre roca total es posible; el potasio--argón arriesgaría edades rebajadas por el argón perdido en un calentamiento posterior. (d)~Radiocarbono sobre el carbonato de la concha, cerca del límite del método (unos \num{70000} años). (e)~Ninguno directamente: su edad queda acotada por rocas volcánicas datadas encima y debajo, o se obtiene correlacionando sus fósiles con una sección datada (\cref{ch:b1:stratigraphic-scale}).
\end{solution}

\begin{solution}{exo:b1:radiometric-methods:3}
(a)~$t = -8033 \ln 0.25 = 8033 \times 1.3863 = \num{11136}$ años BP.
(b)~La vida media es $5700/\ln 2 = \num{8223}$ años: $t = 8223 \times 1.3863 =
\num{11400}$ años. (c)~Ambos suponen que la proporción de carbono-14 en el carbono vivo era entonces la misma que en el estándar; no lo era. La edad de calendario se lee
en una \omterm{def:b1:radiometric-methods:radiocarbon}{curva de calibración} como IntCal20, como un intervalo de años cal BP.
\end{solution}

\begin{solution}{exo:b1:radiometric-methods:4}
(a)~Cada reloj empieza cuando su \omterm{def:b1:rock-forming-minerals:mineral}{mineral} se enfría por debajo de su \omterm{def:b1:radiometric-methods:closure}{temperatura de cierre}:
el circón (por encima de \qty{750}{\celsius}) cierra primero, la hornblenda (unos
\qty{530}{\celsius}) después, y la biotita (unos \qty{280}{\celsius}) la última, así que las edades disminuyen en ese orden a medida que se enfrió el granito. (b)~$(530 - 280)/(410 - 398) =
250/12 = \qty{21}{\celsius}$ por millón de años. (c)~El circón cerró por encima de
\qty{750}{\celsius}, así que el granito se enfrió al menos $750 - 530 =
\qty{220}{\celsius}$ en \qty{10}{\Myr}: al menos \qty{22}{\celsius} por millón de
años, una cota inferior.
\end{solution}

\begin{solution}{exo:b1:radiometric-methods:5}
(a)~$\bar x = 7.5$, $\bar y = 0.809525$, $S_{xx} = 81.0$, $S_{xy} = 1.13965$;
$m = 1.13965/81.0 = 0.0140698$, $b = 0.809525 - 7.5 \times 0.0140698 = 0.70400$.
(b)~$t = \ln(1.0140698)/\num{1.3972e-11} = \qty{1000.0}{\Ma}$; \omterm{def:b1:radiometric-methods:isochron}{cociente inicial}
$0.7040$. (c)~Los residuos son $+1.6$, $-2.5$, $+0.7$ y $+0.2$ (en unidades de
\num{e-4}): $s = \sqrt{\num{9.34e-8}/2} = 0.00022$, $\sigma_m = 0.00022/\sqrt{81}
= \num{2.40e-5}$ y $\sigma_t = \num{2.40e-5}/(\num{1.3972e-11} \times 1.0141) =
\qty{1.7}{\Myr}$: \qty[uncertainty-mode=separate]{1000.0(17)}{\Ma}. (d)~$0.7040$
está cerca del valor del \omterm{def:b1:layered-earth:layers}{manto}, $0.7025$, y bien por debajo del de la \omterm{def:b1:layered-earth:layers}{corteza}, $0.710$: una fuente sobre todo en el \omterm{def:b1:layered-earth:layers}{manto}.
\end{solution}

\begin{solution}{exo:b1:radiometric-methods:6}
(a)~$0.100/39.098 \times \num{1.17e-4} = \qty{2.992e-7}{mol}$ de potasio-40 por
gramo. (b)~${}^{40}\mathrm{Ar}^*/{}^{40}\mathrm{K} = \num{1.169e-9}/\num{2.992e-7}
= 0.0039065$; $\lambda/\lambda_e = 5.543/0.581 = 9.5404$; $t = \ln(1 + 9.5404
\times 0.0039065)/\num{5.543e-10} = \ln(1.037269)/\num{5.543e-10} =
\qty{66.01}{\Ma}$. (c)~El argón radiogénico sería $0.995 \times 1.169$
nanomoles, dando \qty{65.69}{\Ma}: la edad medida sería \qty{0.32}{\Myr}
demasiado antigua, diez veces los \num{32000} años en que el impacto y la
extinción al final del Cretácico se muestran hoy coincidentes. El argón en exceso preocupa sobre todo en rocas jóvenes y en rocas totales; un
\omterm{def:b1:rock-forming-minerals:mineral}{mineral} como la sanidina no debe incorporar argón al cristalizar, y un espectro de edades argón--argón lo revelaría.
\end{solution}

\begin{solution}{exo:b1:radiometric-methods:7}
(a)~Hacia \num{2000} cal BP: la curva pasa por \num{2070} años BP allí.
(b)~De unos \num{2495} a \num{2699} cal BP, doscientos años: la curva es
casi plana entre \num{2500} y \num{2700} cal BP, de modo que cada edad de calendario
en ese intervalo tiene casi la misma \omterm{def:b1:radiometric-methods:radiocarbon}{edad radiocarbónica}, y una medida más precisa no estrecharía el intervalo. (c)~A \num{2000} cal BP la \omterm{def:b1:radiometric-methods:radiocarbon}{edad radiocarbónica} es \num{70} años \emph{más antigua} que la edad de calendario; a \num{8000} cal BP,
\num{822} años más jóvenes (\qty{10}{\percent}); a \num{12000} cal BP, \num{1694}
años más jóvenes (\qty{14}{\percent}). La diferencia no es proporcional a la
edad: incluso cambia de signo. No puede proceder de la semivida sola, que daría un \qty{2.4}{\percent} constante.
\end{solution}

\begin{solution}{exo:b1:radiometric-methods:8}
(a)~La capa de ceniza inferior (\qty{94.50}{\Ma}) da una edad máxima, y la colada de lava superior
(\qty{93.80}{\Ma}) una mínima: el límite queda acotado entre ellas. El dique es también una mínima, pero más débil, y más joven que la lava. El circón detrítico es una máxima solo para la arenisca, y más débil que la capa de ceniza, que yace bajo la arenisca.
(b)~$f = (12 - 4)/(20 - 4) = 0.5$; $t = 94.50 - 0.5 \times 0.70 =
\qty{94.15}{\Ma}$; $\sigma_t = \sqrt{0.25 \times 0.10^2 + 0.25 \times 0.12^2} =
\qty{0.08}{\Myr}$: \qty[uncertainty-mode=separate]{94.15(8)}{\Ma}. (c)~Sus
granos proceden de una roca de al menos \qty{95.3}{\Ma}, erosionada en algún lugar de la zona fuente; muestra que la arenisca es más joven que \qty{95.3}{\Ma}, lo que la capa de ceniza inferior ya demuestra con más ajuste.
\end{solution}

\begin{solution}{exo:b1:radiometric-methods:9}
(a)~Grano~A: $\ln(1.08401)/\num{1.55125e-10} = \qty{520.0}{\Ma}$ y
$\ln(1.6688)/\num{9.8485e-10} = \qty{520.0}{\Ma}$. Grano~B: \qty{368.4}{\Ma} y
\qty{389.9}{\Ma}. (b)~El grano~A es concordante. (c)~$0.05881/0.08401 = 0.7000$ y
$0.4682/0.6688 = 0.7001$. La edad es $\ln(1 + R)/\lambda$, no proporcional a
$R$: para el cociente pequeño del uranio-238, $\ln(1+R) \approx R$ y la edad baja casi un \qty{30}{\percent} (hasta \qty{368}{\Ma}); para el cociente mayor del uranio-235 el logaritmo baja menos, y la edad solo hasta \qty{390}{\Ma}. Ambas son demasiado jóvenes, y discrepan: el grano es discordante.
\end{solution}

\begin{solution}{exo:b1:radiometric-methods:10}
(a)~Con $u = \lambda t$, $D^*/P = e^u - 1 = u + u^2/2 + \dots$, así que $(D^*/P)/\lambda
= t\,(1 + u/2 + \dots)$: la aproximación sobrestima la edad con un error relativo cercano a $u/2$. (b)~Rubidio-87 a \qty{310}{\Ma}: $u = 0.00433$, error
\qty{0.22}{\percent} (\qty{0.7}{\Myr}); potasio-40 a \qty{296}{\Ma}: $u =
0.164$, error \qty{8.7}{\percent} (\qty{26}{\Myr}; $u/2$ da \qty{8.2}{\percent});
uranio-238 a \qty{4.4}{\Ga}: $u = 0.683$, error \qty{43}{\percent} --- la
aproximación e incluso su estimación del error ($u/2 = \qty{34}{\percent}$) fallan.
(c)~$R = 0.004341$, $\sigma_R = \num{4.34e-6}$; por
\cref{prop:b1:radiometric-methods:precision}, $\sigma_t =
\num{4.34e-6}/(\num{1.3972e-11} \times 1.0043) = \qty{0.31}{\Myr}$.
\end{solution}

\begin{solution}{exo:b1:radiometric-methods:11}
(a)~Sea ${}^{40}\mathrm{K}$ el potasio-40 de hoy. En el tiempo $t_2$ antes del
presente había ${}^{40}\mathrm{K}\,e^{\lambda t_2}$ átomos, y el argón producido entre $t_1$ y $t_2$ fue, por \cref{prop:b1:radiometric-methods:kar} aplicada a ese intervalo,
$(\lambda_e/\lambda)\,{}^{40}\mathrm{K}\,e^{\lambda
t_2}\bigl(e^{\lambda (t_1 - t_2)} - 1\bigr) = (\lambda_e/\lambda)\,
{}^{40}\mathrm{K}\,\bigl(e^{\lambda t_1} - e^{\lambda t_2}\bigr)$. Una fracción
$1 - f$ de ella permaneció. Desde $t_2$, el \omterm{def:b1:rock-forming-minerals:mineral}{mineral} cerrado ha producido
$(\lambda_e/\lambda)\,{}^{40}\mathrm{K}\,(e^{\lambda t_2} - 1)$ más. Sumando y
multiplicando por $\lambda/(\lambda_e\,{}^{40}\mathrm{K})$ se obtiene la expresión.
(b)~La edad aparente es $\ln(1 + \text{expresión})/\lambda$: \qty{400.0}{\Myr} para
$f = 0$, \qty{342.3}{\Myr} para $0.3$, \qty{282.7}{\Myr} para $0.6$ y
\qty{200.0}{\Myr} para $1$. (c)~Cada edad entre \qty{200}{\Myr} y
\qty{400}{\Myr} puede producirse por alguna pérdida parcial, y un solo número no puede
decir si registra un suceso real en ese momento o una mezcla; hacen falta un espectro de edades, u otros relojes de la misma roca.
\end{solution}

\begin{solution}{exo:b1:radiometric-methods:12}
(a)~La edad es proporcional a la vida media usada: $10306 \times 8223/8033 =
\num{10550}$ años. (b)~$10550 - 10306 = 244$ años, una séptima parte de los
\num{1694}. (c)~Si la actividad inicial era $k A_0$, la edad verdadera es $t = \tau
\ln(k A_0/A) = \num{10550} + \tau \ln k$; con $t = \num{12000}$, $\ln k =
1450/8223 = 0.176$ y $k = 1.19$: unos \qty{19}{\percent} más carbono-14 que
en el estándar.
\end{solution}

\begin{solution}{pb:b1:radiometric-methods:1}
\begin{enumerate}
  \item Los puntos suben de forma constante con $x$; sus desviaciones de una recta son como mucho
        dos unidades en la cuarta cifra decimal, frente a un intervalo de $0.04$ en $y$: están sobre una recta, con poca dispersión.
  \item $\bar x = 4.90333$, $\bar y = 0.729288$, $S_{xx} = 62.1253$, $S_{xy} =
        0.269335$; $m = 0.269335/62.1253 = 0.0043354$.
  \item $b = 0.729288 - 0.0043354 \times 4.90333 = 0.70803$: la
        $^{87}\mathrm{Sr}/^{86}\mathrm{Sr}$ del granito cuando cristalizó,
        compartida por todas sus partes.
  \item $t = \ln(1.0043354)/\num{1.3972e-11} = \qty{309.6}{\Ma}$.
  \item Los residuos son $+9.1$, $-17.2$, $+6.9$, $-6.4$, $+15.1$ y $-7.5$ (en
        unidades de \num{e-5}); $s = 0.000137$, $\sigma_m = 0.000137/\sqrt{62.1253} =
        \num{1.74e-5}$, $\sigma_t = \num{1.74e-5}/(\num{1.3972e-11} \times 1.0043)
        = \qty{1.2}{\Myr}$: \qty[uncertainty-mode=separate]{309.6(12)}{\Ma}, y
        $\sigma_b = 0.00010$.
  \item $0.7080$ está más cerca del $0.710$ de la \omterm{def:b1:layered-earth:layers}{corteza} que del $0.7025$ del \omterm{def:b1:layered-earth:layers}{manto}.
        Para una mezcla, $(0.70803 - 0.7025)/(0.710 - 0.7025) = 0.74$: unas
        tres cuartas partes del estroncio procedían de la \omterm{def:b1:layered-earth:layers}{corteza} fundida.
  \item $0.0740/39.098 \times \num{1.17e-4} = \qty{2.214e-7}{mol}$ por gramo,
        por $\num{6.022e23}$: \num{1.33e17} átomos.
  \item $^{40}\mathrm{Ar}^*/^{40}\mathrm{K} = \num{4.139e-9}/\num{2.2144e-7} =
        0.018691$; $1 + 9.5404 \times 0.018691 = 1.17832$; $t =
        \ln(1.17832)/\num{5.543e-10} = \qty{296.0}{\Ma}$.
  \item $^{40}\mathrm{K} = \qty{3.292e-8}{mol.g^{-1}}$, cociente $0.019239$,
        $\ln(1.18355)/\num{5.543e-10} = \qty{304.0}{\Ma}$.
  \item $\ln(1.018691)/\num{5.543e-10} = \qty{33.4}{\Ma}$, nueve veces demasiado joven,
        porque solo una desintegración de cada diez da argón. El calcio-40 es el isótopo común del calcio, abundante en las rocas: lo poco producido por desintegración no puede distinguirse del calcio que los \omterm{def:b1:rock-forming-minerals:mineral}{minerales} contenían desde el principio.
  \item El exceso de argón desplaza las edades en cantidades del orden de un millón de años, lo que importa en rocas jóvenes; aquí sería muy por debajo del \qty{1}{\percent} de la edad. Y los cristales de \omterm{def:b1:rock-forming-minerals:groups}{mica} y de \omterm{def:b1:rock-forming-minerals:groups}{anfíbol} apenas incorporan argón al formarse.
  \item O bien el granito se enfrió despacio y las \omterm{def:b1:rock-forming-minerals:groups}{micas} y el \omterm{def:b1:rock-forming-minerals:groups}{anfíbol} cerraron después de cristalizar (\omterm{def:b1:radiometric-methods:closure}{edades de enfriamiento}), o un calentamiento posterior expulsó parte de su argón. Una meseta plana en un espectro de edades argón--argón respalda una \omterm{def:b1:radiometric-methods:closure}{edad de enfriamiento}; un espectro sin meseta advierte de que los \omterm{def:b1:rock-forming-minerals:mineral}{minerales} perdieron argón. El orden de las edades, con la hornblenda antes que la biotita como predicen sus \omterm{def:b1:radiometric-methods:closure}{temperaturas de cierre}, también apunta al enfriamiento.
  \item $\ln(1.049361)/\num{1.55125e-10} = \qty{310.6}{\Ma}$ y
        $\ln(1.35784)/\num{9.8485e-10} = \qty{310.6}{\Ma}$: concordantes.
  \item Circón \qty{310.6}{\Ma}, cierre por encima de \qty{750}{\celsius}:
        cristalización; la \omterm{def:b1:radiometric-methods:isochron}{isócrona} de roca total,
        \qty[uncertainty-mode=separate]{309.6(12)}{\Ma}, también la data,
        y concuerda dentro de su incertidumbre; hornblenda \qty{304.0}{\Ma}: enfriamiento
        por debajo de unos \qty{530}{\celsius}; biotita \qty{296.0}{\Ma}: enfriamiento
        por debajo de unos \qty{280}{\celsius}. La gráfica desciende desde más de
        \qty{750}{\celsius} a \qty{310.6}{\Ma} hasta \qty{280}{\celsius} a
        \qty{296.0}{\Ma}.
  \item $(530 - 280)/(304.02 - 296.03) = 250/7.99 = \qty{31}{\celsius}$ por
        millón de años.
  \item $(750 - 530)/(310.60 - 304.02) = 220/6.58 = \qty{33}{\celsius}$ por
        millón de años como mínimo, ya que el circón cerró por encima de
        \qty{750}{\celsius}: el enfriamiento fue al menos tan rápido al principio.
  \item La diferencia de temperaturas, $250 \pm \sqrt{40^2 + 40^2} = 250 \pm
        57$, da \qtyrange{24}{38}{\celsius} por millón de años.
  \item $(550 - 350)/7.99 = \qty{25}{\celsius}$ por millón de años, dentro del intervalo
        de la pregunta~17: la incertidumbre procede de las \omterm{def:b1:radiometric-methods:closure}{temperaturas de cierre}, no de las edades.
  \item $(280 - 10)/30 = \qty{9.0}{km}$.
  \item $\ln(1.039108)/\num{1.55125e-10} = \qty{247.30}{\Ma}$, y lo mismo a partir del
        cociente del uranio-235: concordantes.
  \item Del granito, erosionado en la \omterm{def:b1:cross-sections:unconformity}{discordancia}: su edad es la de su cristalización. Solo muestran que la arenisca es más joven que
        \qty{310}{\Ma}: una edad máxima, mucho más débil que la capa de ceniza.
  \item $(58 - 18)/(247.30 - 245.30) = \qty{20.0}{m}$ por millón de años; la base
        está $18/20.0 = \qty{0.90}{\Myr}$ más antigua que la ceniza inferior:
        \qty{248.20}{\Ma}.
  \item Desde \qty{309.6}{\Ma} (\omterm{def:b1:radiometric-methods:isochron}{isócrona}) o \qty{310.6}{\Ma} (circón) hasta
        \qty{248.2}{\Ma}: unos \qty{62}{\Myr}.
  \item \qty{9.0}{km} en $296.0 - 248.2 = \qty{47.8}{\Myr}$: \qty{0.19}{km} por
        millón de años, \qty{0.19}{mm} por año. Esto supone que la biotita datada procedía justo por debajo de la \omterm{def:b1:cross-sections:unconformity}{discordancia}, de modo que toda la roca que la cubría en el granito y en el esquisto se eliminó cuando llegó la arena, y que el \omterm{def:b1:heat-budget:geotherm}{geotermo} se mantuvo constante.
  \item El granito se enfrió a unos \qty{31}{\celsius} por millón de años entre
        los cierres de la hornblenda y de la biotita.
\end{enumerate}
\end{solution}
